% Transcription — fonds Alexandre Grothendieck, shelfmark 85, batch 4 (pages 61-80).
% Established from the facsimile archives/batches/85/batch-04.pdf (a copy of
% 85.pdf fetched through the project's relay,
% https://grothendieck-relay.onrender.com/source/85.pdf), per the skill
% transcribe-grothendieck. The folder is ninety-seven pages in five batches;
% this is the fourth.
%
% Pass: Opus 5.5 (claude-opus-5-5), 2026-09-26 — first pass, unchecked against the pages by a human.
%
% Pages skipped: none as unrelated. Page 74 is a cover on pink card,
% inscribed in ink at the top right « 2-Polyèdres réguliers finis en car. 0
% (classification) » and in pencil « (6 p) », which counts exactly pages
% 75-80; it takes no \page marker and its words become the section heading,
% with a \note. Pages transcribed: 61-73, 75-80 (nineteen). Pages
% summarised: none. On p. 64 the lower half carries, head to tail and
% cancelled by oblique strokes, an earlier draft (three pencils of lines,
% formulas in rho, eta, lambda_0 = phi_0^2); it is given in one \note with its
% legible formulas. Pages 71 and 79 are crossed through by long strokes and
% are transcribed as read, the cancellation recorded once.
%
% His pagination: none of his own on these leaves. Pages 61-73 continue the
% run « Hyperpolygônes réguliers » opened by the cover of p. 57 (batch 3),
% whose pencil count « (16 p) » covers pages 58-73. Numbered and lettered
% steps: a)-b) « Principes d'élimination » (p. 62); « En résumé » a)-b)
% (p. 65); 1)-2) (p. 66); the lettered conditions a), b), b'), b'') (p. 75);
% Prop. a)-c), the Theorem's cases I, II, III, II', III', IV (p. 76), over a
% cancelled first list 1), 2), 2'), 3), 3'); « Indications sur la dém. du
% th. » I), II), III) (pp. 77-79). No date in his hand anywhere in the batch.
%
% What the batch is about. Pages 61-73 end the run on regular hyperpolygons
% in characteristic 0 (case n = 1, two strict pseudo-reflections tau_0, tau_1
% of a plane V with eigenvalues lambda_i = exp(i theta_i), invariant beta =
% xi (lambda_0 - 1)(lambda_1 - 1)). P. 61: xi satisfies a quadratic equation
% (E) whose constant term involves kappa = (Tr u)^2 / det u - 2 for u =
% tau_0 tau_1, kappa in {-2, -1, 0, alpha, alpha'}; the discriminant is
% (kappa + 2) / (4 sin^2 sin^2) >= 0, zero only when tau_0^C tau_1^C has
% order 2. Pp. 62-65: principles of elimination (0 < xi < 1; stability under
% theta_0 -> i theta_0, theta_1 -> j theta_1); the projective situation is
% determined by kappa_0, kappa_1, kappa (the Schwartz system), four
% hyperpolygonal lifts; a table of correspondences; « En résumé »: for
% nu_0, nu_1 != 2 one hyperpolygon per admissible (lambda_0, lambda_1,
% kappa), for nu_0 or nu_1 = 2 two. Pp. 66-68: the list of numerical types
% giving two hyperpolygons, the table of 6 + 13 + 25 = 44 cases with the
% number of hyperpolygons in each (total 4.45 = 180). Pp. 69-73: the test
% f_kappa(1) >= 0 (both roots of (E) below 1), reduced to cos((theta_1 -
% theta_0)/2) + cos(Theta/2) >= 2 sin(theta_0/2) sin(theta_1/2), checked
% case by case with alpha^2 + alpha - 1 = 0: (3,3,3), (4,4,3), (4,3,4),
% (3,4,4), (5,5,3), (5,3,5), (3,5,5), (5,5,5) OK; (5,3,5') and (3,5,5')
% left undetermined (« pas vrai ! »). Pages 75-80 (« 2-Polyèdres réguliers
% finis en car. 0 »): case n = 2 with invariants beta_i = alpha_i + 2, rho,
% sigma, sigma'; the conic Q_C = Y and the surface Q_X, smooth when ...;
% the category of pinned regular 2-polyhedra as triples (Y, phi) with
% conditions a), b), b'), b''); Prop. on finiteness (orders nu_i of the
% eigenvalues zeta_i); Theorem: G_0 finite iff one of I (tetrahedron), II
% (octahedron), III (icosahedron), II' (cube), III' (dodecahedron), IV
% (pentagonal, nu_0 = nu_1 = 5, alpha_0 != alpha_1), 3 + 3x2 = 9 types over
% a field containing alpha; sketch of proof via Hurwitz (ramification
% (2,3,3), (2,3,4), (2,3,5)), the standard Pythagorean groups A4, S4 x Z/2,
% A5 x Z/2, the fundamental group of the sphere minus three points, the
% bound 3 <= nu_0, nu_1 <= 5, the positive-definite quadratic form, and the
% elimination of (4,4) and (4,5): « On a gagné ! ».
%
% Notation fixed once, aligned on batch 3: \mathfrak{S}, \mathfrak{A} for
% the symmetric and alternating groups (his two tall letters are close; on
% p. 78 the first is flagged); \check{\mathbf{P}}(V); tau_i^C for the image
% in Aut(C); alpha, alpha' for the two roots of alpha^2 + alpha - 1 = 0 and
% 5, 5' in the types; kappa for his K (trace invariant). His « OPS » is kept.
% Underlined words are \emph; his boxes are \boxed in mathematics and noted
% in prose.
%
% Hardest pages: 61 (dense, struck calculations, a rotated margin column),
% 63 and 65 (fast prose), 67 (the table and its NB), 71 and 72 (numerical
% work crossed through and overwritten), 75 (margin notes written slantwise),
% 77 (a boxed cancelled block). Readings to check first: the equation of
% p. 61 (factor before the bracket), the count « 16 » / « 15 » at its foot,
% the sign in lambda_i^2 + kappa_i lambda_i + 1 = 0 (p. 64), the subscript
% counts of p. 67, « cos Theta = -1 » on p. 71, and the groups G_0 on p. 78.

\documentclass[11pt,a4paper]{article}
\input{../preamble/grothendieck}

\folder{85}
\batch{4}
\pages{61}{80}
\dating{s.d. — le groupe « Géométrie et topologie combinatoire » (69 à 102) est daté 1976-[vers 1986]}
\watermark{Édition de démonstration}
\foldertitle{2-polyèdres et 1-polygônes : notes manuscrites (s.d.)}

\begin{document}

\section*{Hyperpolygônes réguliers (suite)}
\note{Les pages 61 à 73 continuent la suite ouverte par la couverture de la
page 57 (lot 3), «~Hyperpolygônes réguliers~», dont le compte au crayon
«~(16 p)~» couvre les pages 58 à 73. La page 61 reprend au milieu du
calcul.}

\page{61}
\[
\beta = c\,\frac{\cdots}{\cdots}
= c\,\frac{(\lambda_0 - 1)(\lambda_1 - 1)}{(\kappa_1 - 1)}
= \xi\,(\lambda_0 - 1)(\lambda_1 - 1) = \xi\,\mu_0\mu_1,
\qquad \xi = \frac{c}{\kappa_1 - 1} \in \mathbb{R}
\]
\[
\kappa_i = \mathrm{Tr}_{\mathbb{C}/\mathbb{R}}\,\lambda_i
= \lambda_i + \overline{\lambda}_i = \lambda_i + \lambda_i^{-1}
\qquad (\lambda_i \overline{\lambda}_i = 1)
\]
\emph{NB} $\det \tau_0\tau_1 = \lambda_0\lambda_1$,
$\mathrm{Tr}\,\tau_0\tau_1 = $
\note{La première égalité est surchargée ; le quotient qui suit $c$, noté
ici par des points, ne se lit pas (\ill{}).}

$\xi$ satisfait à l'équation
\[
\frac{(\lambda_0 - 1)(\lambda_1 - 1)}{\lambda_0\lambda_1}
\Bigl[\xi^2 + 2\,\frac{\lambda_0 + \lambda_1}{(\lambda_0 - 1)(\lambda_1 - 1)}\,\xi
+ \frac{\lambda_0^2 + \lambda_1^2}{(\lambda_0 - 1)^2(\lambda_1 - 1)^2}\Bigr]
= \kappa \in \lbrace -2, -1, 0, \alpha, \alpha' \rbrace
\]
qui exprime que $\tau_0^{C}\tau_1^{C}$ est d'ordre 2, 3, 4 ou 5.
[$u = \tau_0\tau_1$ a comme trace $\lambda_0 + \lambda_1 + \beta$, comme
déterminant $\lambda_0\lambda_1$], et
\[
\kappa = \kappa(u) = \frac{(\mathrm{Tr}\, u)^2}{\det u} - 2
\]
\note{Le facteur devant le crochet est écrit dans la marge gauche ; sa
lecture et celle du dernier terme sont incertaines.}

\emph{NB} si $\lambda = \exp i\theta$, $\lambda' = \exp i\theta'$, alors
\[
|\lambda + \lambda'| = 2\Bigl|\cos\frac{\theta' - \theta}{2}\Bigr|,\quad
|\lambda^2 + \lambda'^2| = 2|\cos(\theta' - \theta)|,
\]
\[
|\lambda - 1|^2 = 2 - \kappa(\lambda),\quad
\kappa(\lambda) = \mathrm{Tr}_{\mathbb{C}/\mathbb{R}}\,\lambda = 2\cos\theta,\quad
|\lambda - 1| = 2\Bigl|\sin\frac{\theta}{2}\Bigr|
\]
\[
2\,\frac{\lambda + \lambda'}{(\lambda - 1)(\lambda' - 1)}
= -\,\frac{\cos\frac{\theta' - \theta}{2}}
{\sin\frac{\theta}{2}\sin\frac{\theta'}{2}},
\qquad
\frac{(\lambda_0 - 1)^2(\lambda_1 - 1)^2}{\lambda_0\lambda_1}
= 16\sin^2\frac{\theta_0}{2}\sin^2\frac{\theta_1}{2}
\]
\[
\frac{\lambda^2 + \lambda'^2}{(\lambda - 1)^2(\lambda' - 1)^2}
= \frac{2\cos(\theta' - \theta)}{16\sin^2\frac{\theta}{2}\sin^2\frac{\theta'}{2}}
= \frac{1}{8}\,\frac{\cos(\theta' - \theta)}{\sin^2\frac{\theta}{2}\sin^2\frac{\theta'}{2}},
\qquad 0 < \theta_0, \theta_1 \leq \pi
\]
d'où
\[
\xi^2 - \frac{\cos\frac{\theta_1 - \theta_0}{2}}{\sin\frac{\theta_0}{2}\sin\frac{\theta_1}{2}}\,\xi
+ \frac{1}{8}\Bigl(\frac{2\cos(\theta_1 - \theta_0) - \kappa}
{\sin^2\frac{\theta_0}{2}\sin^2\frac{\theta_1}{2}}\Bigr) = 0
\]
\[
\Delta(\kappa) = \frac{1}{2}\,\frac{\cdots}{\sin^2\frac{\theta_0}{2}\sin^2\frac{\theta_1}{2}}
\Bigl[\underbrace{2\cos^2\frac{\theta_1 - \theta_0}{2} - \cos(\theta_1 - \theta_0)}_{1}
+ \frac{1}{2}\kappa\Bigr]
= \boxed{\frac{\kappa + 2}{4\sin^2\frac{\theta_0}{2}\sin^2\frac{\theta_1}{2}}} \geq 0
\]
\[
\kappa + 2 = 2(1 + \cos\Theta) = 4\cos^2\frac{\Theta}{2}
\]
\note{Dans la première fraction, un 4 biffé au numérateur et au
dénominateur ; dans l'équation en $\xi$, un terme biffé après $\kappa$ ;
dans $\Delta(\kappa)$, le numérateur (points) est biffé et illisible
(\struck{\ill{}}). Plusieurs états successifs sont biffés ; l'expression encadrée est entourée d'un cercle. Une colonne de
calculs écrite en travers dans la marge gauche (rapports de sinus et de
cosinus des demi-angles) ne se lit qu'en partie.}

\struck{\ill{} \ill{} \ill{}} N'est nul que si $\kappa = -2$ i.e.\
$\tau_0^{C}\tau_1^{C}$ est d'ordre 2. \struck{\ill{}} auquel cas il y a une
seule solution
\[
\boxed{\xi = \frac{\cos\frac{\theta_1 - \theta_0}{2}}{2\sin\frac{\theta_0}{2}\sin\frac{\theta_1}{2}}}
\]
\ill{} les valeurs de $\lambda_i = \exp 2i\pi/\nu_i$, $2 \leq \nu_i \leq 5$.
Cela fait \uncertain{16} \add{$= 15$} possibilités (en excluant
$\nu_0 = \nu_1 = 2$).
\note{Le «~16~» est repassé en gras et barré d'un trait oblique, avec
«~$= 15$~» au-dessus ; la lecture du chiffre retenu est incertaine.}

\page{62}
Pour chacun des \add{quatre autres} cas $\kappa = -1, 0, \alpha, \alpha'$,
cela fait deux autres possibilités, ce qui fait a priori encore
$8 \cdot 15 = 120$ possibilités --- mais il faut en exclure un bon
paquet\ldots~!

Principes d'élimination
\begin{itemize}
\item[a)] la condition
\[
0 < \xi < 1
\]
exprimant que la forme \struck{\ill{}} \ill{} est définie
\uncertain{(ou $=$)},
\item[b)] la condition que $\xi$ satisfasse des équations \ill{} quand on
change $\theta_0$ en $i\theta_0$ ($i \not\equiv 0\ (\nu_0)$) et $\theta_1$
en $j\theta_1$ ($j \not\equiv 0\ (\nu_1)$), quitte à changer $\kappa$ dans
$\lbrace -2, -1, 0, \alpha, \alpha' \rbrace$ --- i.e.\ que \ill{} des
valeurs prises par
\[
16\sin^2\frac{\theta}{2}\sin^2\frac{\theta'}{2}
\Bigl(\xi^2 - \frac{\cos\frac{\theta' - \theta}{2}}{2\sin\frac{\theta}{2}\sin\frac{\theta'}{2}}\,\xi
+ \frac{1}{8}\,\frac{\cos(\theta' - \theta)}{\sin^2\frac{\theta}{2}\sin^2\frac{\theta'}{2}}\Bigr)
\]
pour $\theta = i\theta_0$, $\theta' = j\theta_1$ ($i \not\equiv 0\ (\nu_0)$,
$\theta' \not\equiv 0\ (\nu_1)$) --- soit dans
$\lbrace -2, -1, 0, \alpha, \alpha' \rbrace$\ldots
\end{itemize}
\note{Au dénominateur du coefficient de $\xi$, les demi-angles sont
surchargés.}

L'analyse géométrique des «~syst.\ de Schwartz~» montre que la
\emph{situation projective} (dans $C = \check{\mathbf{P}}(V)$) est
déterminée mod iso, par la connaissance de $\kappa_0$, $\kappa_1$ et
$\kappa$ \struck{(qui fixent \ill{} $\kappa_0, \kappa_1$)} (classes de
conjugaison de $\tau_0^{C}, \tau_1^{C}$ et $(\tau_0^{C}\tau_1^{C})^{-1}$ ou
$\tau_0^{C}\tau_1^{C}$) dans \struck{D'autant plus} $G = \mathrm{Aut}(X)$,
\ill{} ce qui revient au même, dans le groupe pythagoricien $H$ qu'ils
engendrent (mod $\mathrm{Aut}(H)$ dans le cas tétraédral, \uncertain{mod}
$\mathrm{Int}(H)$).

\page{63}
D'autre part, pour une situation projective donnée, il y a 4 choix de
situations hyperpolygônales i.e.\ donnant naissance (savoir le choix d'un
des deux pts fixes pour $\tau_0$, et pour $\tau_1$) au prix d'une telle
situation. \uncertain{Correspondance}~: $\lambda_0, \lambda_1$,
$\beta = \xi(\lambda_0 - 1)(\lambda_1 - 1)$. D'un autre côté, remplaçant
$\tau_0$ par $\frac{1}{\lambda_0}\tau_0$, ou $\tau_1$ par
$\frac{1}{\lambda_1}\tau_1$, ou l'un et l'autre à la fois --- ce qui ne
change pas $\kappa_0, \kappa_1, \kappa'$ car $\tau_0, \tau_1$ et
$\tau_0\tau_1$ sont multipliés chacun par un scalaire. \struck{Donc en
principe} Mais ceci remplace $(\lambda_0, \lambda_1)$ par
$(\lambda_0^{-1}, \lambda_1)$ resp.\ $(\lambda_0, \lambda_1^{-1})$ resp.\
$(\lambda_0^{-1}, \lambda_1^{-1})$ qui sont \struck{déjà différents de}
\add{«~en général~»} \ill{} \ill{} des valeurs \ill{}
\struck{\ill{}}~:

\struck{\ill{} des valeurs correspondantes de $\beta$, qui ne peut
\struck{\ill{}} \add{a priori} prendre que 2 val.\ possibles, d'après
l'équation \ill{} plus \ill{} $\beta$}
\note{Bloc encadré et barré de deux traits obliques.}

--- «~en général~» signifiant que $\nu_0$ et $\nu_1$ sont $\neq 2$, i.e.\
$\lambda_0, \lambda_1 \neq -1$ (i.e.\ $\mu_0, \mu_1 \neq -2$). Dans ce cas il
s'ensuit \add{$\lambda_0, \lambda_1$ et $\kappa$ fixés,
\struck{\ill{}} $\kappa \neq -2$} \struck{les 2 solutions
\ill{}} de ce que \struck{\ill{}} l'équation en $\xi$ a a priori 2 sol.
\note{La fin de la page est blanche.}

\page{64}
\emph{Groupes pythagoriciens}

$\updownarrow$

\struck{\ill{}} \emph{données}
$(\lambda_0, \lambda_1, \beta = \xi(\lambda_0 - 1)(\lambda_1 - 1))$

$\updownarrow$

\struck{\ill{}} couples de 2 ps.\ réfl.\ $\tau_0, \tau_1$ dans \add{plan}
vectoriel $V$ (épinglés par deux droites $a_0 \in V^{\tau_0}$\ldots)~:
conditions de non dég.\ et d'engendrer d'un groupe fini\ldots
\note{Les trois termes sont reliés par des flèches verticales à double
pointe. Sous le dernier, un petit schéma : $V$, les droites
$L_0, L_1$ et $L_0', L_1'$, reliées par des traits.}

$\longleftrightarrow$ \struck{couples pythag.} \add{\uncertain{couples
pythag.}} (\struck{\ill{}}) d'un $(\tau_0^{C}, \tau_1^{C})$ dans plan proj.\
$+$ $a_0$ point fixe de $\tau_0^{C}$, $a_1$ pt fixe de $\tau_1^{C}$
$\longleftrightarrow$ ceci équivaut à donnée de la situation pyth.\
\emph{projective} $\kappa_0, \kappa_1, \kappa$ admis., $+$ $a_0, a_1$ dans le
syst.\ défini par $\kappa_0, \kappa_1, \kappa$.
\note{Colonne de droite de la page, reliée à la précédente par une flèche
horizontale à double pointe.}

Si $\kappa_i$ fixé, $\lambda_i$ sont déterminés comme \ill{} des deux
solutions de
\[
\lambda_i^2 + \kappa_i\lambda_i + 1 = 0
\]
Le choix des $\lambda_i$ équivaut, \struck{\ill{}} pour $\lambda_i \neq -1$
i.e.\ $\kappa_i \neq -2$ i.e.\ $\nu_i \neq 2$, au choix des $a_i$
\struck{\ill{}}~; si par contre $\nu_i = 2$ i.e.\ $\lambda_i = -1$ i.e.\
$\kappa_i = -2$,
\note{Le signe de $\kappa_i\lambda_i$ est transcrit tel qu'écrit ; avec
$\kappa_i = \lambda_i + \lambda_i^{-1}$ (page 61) on attendrait
$\lambda_i^2 - \kappa_i\lambda_i + 1 = 0$. La phrase s'arrête sur la
virgule.}

\note{La moitié inférieure de la page porte, tête-bêche et barré de traits
obliques, un brouillon antérieur : trois faisceaux de droites concourantes
(légendes $D_1, D_2, D_1'$ et $D, D', D''$), et les formules
«~$\rho(e_0) = e_0$, $\rho(e_1) = -e_1$, $\rho'(e_0) = $~»,
«~$(\varepsilon_1, \varepsilon_2) = \eta$, $\eta^2 = \mathrm{id}$,
$\eta'^2 = \mathrm{id}$, $\eta''^2 = \mathrm{id}$, $\eta\eta' = ?$,
$\eta'\eta'' = ?$, $\eta\eta'\eta'' = $~», et
«~$\lambda_0 = \varphi_0^2 = (s_0 + ic_0)^2 = s_0^2 - c_0^2 + 2is_0c_0$~»,
«~$\lambda_1 = \varphi_1^2 = $~».}

\page{65}
\struck{Donc le choix} (ce qui ne change pas dans le cas \ill{} pour un
seul des indices $i = 0, 1$) alors la donnée de $(\lambda_0, \lambda_1,
\kappa)$ \ill{} détermine \uncertain{pas} la classe d'iso des hyperpolygônes
pythagoriciens, i.e.\ \ill{} déterminée par $\xi$.

\emph{En résumé} (\ill{} \ill{} \uncertain{projectif})
\begin{itemize}
\item[a)] Pour $(\lambda_0, \lambda_1, \kappa)$ admissible (d'où
$\kappa_0, \kappa_1$), si $\kappa_0, \kappa_1 \neq -2$ i.e.\
$\lambda_0, \lambda_1 \neq -1$ i.e.\ $\nu_0, \nu_1 \neq 2$, il y a un
\add{et un} seul hyperpolygône associé (\uncertain{mod.}\ iso) i.e.\ un et
un seul $\xi$ correspondant. Si $\kappa = -2$, \add{i.e.\ $\nu = 2$,} il
\ill{} donné \add{comme l'unique solution de} \struck{\ill{}} l'équation en
$\xi$~; si $\kappa \neq -2$ parmi les 2 solutions de cette équation, une
seule convient pour \add{définir} un hyperpolygône pythagoricien.
[Comment la distinguer~? Si on travaille dans $\mathbb{C}$, on peut songer à
la condition $0 < \xi < 1$\ldots]
\item[b)] Par contre si $\kappa_0$ ou $\kappa_1 = -2$ \add{i.e.\ $\nu_0$ ou
$\nu_1 = 2$} i.e.\ $\lambda_0$ ou $\lambda_1 = -1$, alors dès que
$(\lambda_0, \lambda_1, \kappa)$ admissible \add{\uncertain{projective}}
\ill{}, il y a \emph{deux} hyperpolygônes non iso qui lui correspondent,
correspondant aux 2 solutions distinctes $\xi$ de (E) [\emph{NB} on sait
\add{par ailleurs} \ill{} \ill{} alors $\kappa \neq -2$ i.e.\ que l'on a
bien 2 racines distinctes].
\end{itemize}
\note{Dans b), le signe après $\kappa$ dans le crochet final se lit $=$ ou
$\neq$ ; $\neq$ est retenu parce que la phrase conclut à deux racines
distinctes.}

Il n'est donc pas vrai que $\nu_0, \nu_1, \nu$ déterminent la classe d'iso
d'un hyperpolygône, --- à

\page{66}
un normalisé \add{$\lambda_0, \lambda_1$} par le choix de racines
primitives 3\textsuperscript{e}, 4\textsuperscript{e} et
5\textsuperscript{e} de 1 (par exemple $\exp 2i\pi/\nu$ quand on est sur
$\mathbb{C}$), et ceci pour deux raisons~:

1) Cas où $\lambda_0$ ou $\lambda_1$ est égal à $-1$, i.e.\ $\nu_0$ ou
$\nu_1$ égal 2 --- cela introduit (pour $\kappa$ déjà fixé) une ambiguïté
d'ordre 2~:
\[
\begin{array}{lll}
\text{Cas} & (3,2,3)\,(2,3,3) & (\text{tétraédraux}) \\
 & (3,2,4)\,(2,4,3)\,(2,3,4)\,(4,2,3) & (\text{octaédraux}) \\
 & (3,2,\underline{5})\,(2,5,3)\,(2,3,\underline{5})\,(5,2,3) & (\text{icosaédraux}) \\
 & (5,2,5')\,(2,5,5') &
\end{array}
\]
\note{Les deux dernières lignes sont réunies par une accolade.}

2) Cas où $\nu = 5$, et où on peut avoir le choix entre $\kappa = \alpha$ et
$\kappa = \alpha'$ (cas icosaédral)
\[
\begin{array}{ll}
(3,2,5)\,(2,3,5) & (\text{qui figurent déjà dans liste 1)}) \\
\bigl[(3,2,5')\,(2,3,5')\bigr] &
\end{array}
\]

Finalement, compte tenu du choix des racines primitives
3\textsuperscript{e} 4\textsuperscript{e} et 5\textsuperscript{e} de 1 pour
normaliser $\lambda_0, \lambda_1$, on trouve les possibilités suivantes ---
où on \uncertain{souligne} les possibilités numériques qui donnent lieu à
\emph{deux} hyperpolygônes, en interprétant comme d'habitude 5, 5' comme
signifiant $\kappa = \alpha$ resp.\ $\kappa = \alpha'$~:

\page{67}
\[
\begin{array}{rll}
6 & \textit{Cas tétraédraux} &
(3,3,2)_{(4)}\ \underline{(3,2,3)}_{(2)}\ \underline{(2,3,3)}_{(2)}\
(3,3,3)_{(4)} \\[4pt]
13 & \textit{Cas octaédraux} &
(4,3,2)_{(4)}\ \underline{(3,2,4)}_{(2)}\ \underline{(2,4,3)}_{(2)}\
\underline{(2,3,4)}_{(2)}\ (3,4,2)_{(4)}\ \underline{(4,2,3)}_{(2)} \\
 & & (4,4,3)_{(4)}\ (4,3,4)_{(4)}\ (3,4,4)_{(4)} \\[4pt]
10 & \textit{Cas icosaédraux} &
(5,3,2)_{(8)}\ \underline{(3,2,5)}_{(2)}\ \underline{(2,5,3)}_{(4)}\
\underline{(2,3,5)}_{(2)}\ (3,5,2)_{(8)}\ \underline{(5,2,3)}_{(4)} \\
+4 & & \underline{(3,2,5')}_{(2)}\ \underline{(2,3,5')}_{(2)} \\
+4 & & (5,5,3)_{(16)}\ (5,3,5)_{(8)}\ (3,5,5)_{(8)} \\
2 & & (5,3,5')_{(8)}\ (3,5,5')_{(8)} \\
4 & & \underline{(2,5,5')}_{(4)}\ \underline{(5,2,5')}_{(4)} \\
1 & & (5,5,5)_{(16)} \\
\hline
25 & &
\end{array}
\]
\note{Les nombres entre parenthèses sont écrits sous chaque triple ;
ceux de la ligne «~$+4$~» des cas icosaédraux, et les lignes «~$+4$~» et
«~2~», sont regroupées par des accolades. Dans $(5,3,5')$ et $(3,5,5')$ les
premiers chiffres sont surchargés.}

\[
\begin{array}{rl}
6 & \text{Cas tétraédraux} \\
13 & \text{Cas octaédraux} \\
25 & \text{Cas icosaédraux} \\
\hline
44 & \text{Cas au total}
\end{array}
\]
[Il y a 14 cas sur les 44 où le type numérique ne suffit pas à déterminer
le type d'iso.\ \ill{}]

\emph{NB} On a mis \ill{} \ill{}, en \ill{}, le nombre d'hyperpolygônes
\ill{} \ill{} \ill{} \ill{} \uncertain{Ta-Ti}~: [égal à
$\varphi(\nu_0)\varphi(\nu_1)$, prenant les valeurs 2, 4, 8, 16]. On vérifie
bien que le total (obtenu \ill{} l'un de tous les cas) redonne
$4 \cdot 45 = 180$. Dans le calcul, il faut bien tenir compte aussi des
\uncertain{conjugaisons}, \ill{} \ill{}
\note{Cette note est écrite en colonne, à droite de la page ; la fin se
perd dans des surcharges.}

[\emph{NB} Les syst.\ pythagoriciens de 2 ps.\ réflexions $\tau_0, \tau_1$
(normalisés par $\lambda_0, \lambda_1$) sont plus nombreux, chacun des 45
systèmes de Schwartz $(\kappa_0, \kappa_1, \kappa)$ possibles donnant
$4$ \add{$= 2 \times 2$} \emph{possibilités}, d'où 180 possibilités. Pour
calculer les $\xi$ dans chacun des 180 cas, en trouvant \ill{} qu'il
faudrait \ill{} vérifier que $\xi$ \ill{}

\page{68}
change quand on remplace $\tau_0, \tau_1$ par
$\tau_0^{\ell_0}, \tau_1^{\ell_1}$ (avec
$(\ell_0, \nu_0) = (\ell_1, \nu_1) = 1$) --- cela qui \ill{} bien \ill{}
pas changer. Ceci \ill{} permet de \ill{} ramener aux 45 cas
\emph{normalisés}, \struck{où il y a une formule} qu'il faudrait
expliciter\ldots

\struck{Résumons les ambiguïtés}
\note{Le reste de la page est blanc.}

\page{69}
\[
f_{\kappa}(\xi) = \xi^2 - \frac{\cos\frac{\theta_1 - \theta_0}{2}}
{\sin\frac{\theta_0}{2}\sin\frac{\theta_1}{2}}\,\xi
+ \frac{1}{\cdot}\,
\frac{\cdot\cos(\theta_1 - \theta_0) - \kappa\cos\Theta}
{\sin\frac{\theta_0}{2}\sin\frac{\theta_1}{2}},
\qquad \kappa = 2\cos\Theta
\]
\[
f_{\kappa}(1) = \sin\frac{\theta_0}{2}\sin\frac{\theta_1}{2}
- \sin\frac{\theta_0}{2}\sin\frac{\theta_1}{2}\cos\frac{\theta_1 - \theta_0}{2}
+ \frac{1}{\cdot}\bigl(\cdot\cos(\theta_1 - \theta_0) - \cos\Theta\bigr)
\]
\note{Les chiffres notés $\cdot$, devant $\cos$ et au dénominateur, sont
surchargés et illisibles (\ill{}) ; le
dernier numérateur porte «~$\kappa\cos$~» ou «~$\cos$~» selon l'état.}

\[
\frac{2\pi}{5} \leq \theta_0, \theta_1, \Theta \leq \frac{2\pi}{3}
\]
\[
|\theta_1 - \theta_0| \leq \frac{2\pi}{3} - \frac{2\pi}{5} = \frac{4\pi}{15}
\quad\Bigl(< \frac{2\pi}{5} \leq \Theta, \theta_0, \theta_1\Bigr)
\]
\struck{sauf \ill{} $\Theta = \frac{4\pi}{5}$ i.e.\ $\kappa = \alpha'$}

$\cos(\theta_1 - \theta_0)$ \struck{$\geq$} \ill{} $\cos\Theta$ \ill{} dans
les termes constants de (E) sont $> 0$

\struck{$\cos(\theta_1 - \theta_0) \geq \cos\Theta$, si
$|\theta_1 - \theta_0| \leq \Theta$ p.ex.}\
\struck{les deux racines \ill{} $\Theta \geq \frac{4\pi}{15}$, \ill{}}
\drawing{Sous ces lignes, un croquis barré de courbes oscillantes, et un
schéma de parabole coupant un segment $[0, 1]$ en deux points.}

\[
(\xi - \xi')^2 = (\xi + \xi')^2 - 4\xi\xi'
\]
\[
= \frac{\cos^2\frac{\theta_1 - \theta_0}{2}}
{\sin^2\frac{\theta_0}{2}\sin^2\frac{\theta_1}{2}}
- \frac{1}{2}\,\frac{\cos(\theta_1 - \theta_0) - \cos\Theta}
{\sin^2\frac{\theta_0}{2}\sin^2\frac{\theta_1}{2}} \geq 1\ ?
\]
\[
\underbrace{\cdot\,\sin^2\frac{\theta_0}{2}\sin^2\frac{\theta_1}{2}}_{\geq\ 2\sin^4\frac{\pi}{3}}
+ \bigl(\cos(\theta_1 - \theta_0) - \cos\Theta\bigr)
\]
\note{Le coefficient noté $\cdot$ est illisible (\ill{}).}

\[
7\ \text{cas}\ \left\lbrace
\begin{array}{l}
(3,3,3) \\
(4,4,3)\,(4,3,4)\,(3,4,4) \\
(5,5,3)\,(5,3,5)\,(3,5,5) \\
(3,5,5')\,(5,3,5') \\
(5,5,5)
\end{array}\right.
\]
\note{Dans la liste, $(4,3,4)(3,4,4)$, $(5,3,5)(3,5,5)$ et
$(3,5,5')(5,3,5')$ sont chacun réunis par une accolade.}

\page{70}
À vérifier
\[
\frac{1}{4\sin\frac{\theta_0}{2}\sin\frac{\theta_1}{2}}
\cdot\Bigl[2\cos\frac{\theta_1 - \theta_0}{2} + \sqrt{\kappa + 2}\Bigr] \geq 1
\]
i.e.
\[
2\cos\frac{\theta_1 - \theta_0}{2}
+ \underbrace{\sqrt{\kappa + 2}}_{2\cos\Theta/2}
\geq 4\sin\frac{\theta_0}{2}\sin\frac{\theta_1}{2}
\qquad \bigl(\kappa + 2 = 2(\cos\Theta + 1) = 4\cos^2\tfrac{\Theta}{2}\bigr)
\]
i.e.
\[
\boxed{\cos\frac{\theta_1 - \theta_0}{2} + \cos\frac{\Theta}{2}
\geq 2\sin\frac{\theta_0}{2}\sin\frac{\theta_1}{2}}
\]

\[
\boxed{(3,3,3)}\qquad
\underbrace{1 + \tfrac{1}{2}}_{3/2} \geq
\underbrace{2\Bigl(\frac{\sqrt{3}}{2}\Bigr)^2}_{3/2}\quad \text{OK}
\]
\[
\boxed{(4,4,3)}\qquad
\underbrace{1 + \tfrac{1}{2}}_{3/2} \geq
\underbrace{2\Bigl(\frac{\sqrt{2}}{2}\Bigr)^2}_{1}\quad \text{OK}
\]
\note{Dans la deuxième boîte, le premier chiffre est un 4 écrit sur un 3.}
\[
\boxed{(4,3,4)}\ \text{et}\ \boxed{(3,4,4)}\qquad
\frac{1}{2}\sqrt{2 + \sqrt{3}} + \frac{\sqrt{2}}{2}
\geq \underbrace{2\Bigl(\frac{\sqrt{6}}{4}\Bigr)}_{\sqrt{6}/2}
\]
\[
\text{i.e.}\quad \sqrt{2 + \sqrt{3}} + \sqrt{2} \geq \sqrt{6},\qquad
2 + \sqrt{3} + 2 + \underbrace{2\sqrt{4 + 2\sqrt{3}}}_{\geq 4} \geq 6
\quad \text{OK}
\]
\[
\boxed{(5,5,3)}\qquad
1 + \frac{1}{2} \geq
2\underbrace{\sin^2\frac{\pi}{5}}_{\frac{1}{4}(2 - \alpha)}
= \frac{1}{2}(2 - \alpha)
\]
\[
3 \geq (2 - \alpha),\qquad 1 \geq -\alpha\quad \text{OK}
\]
\note{Dans la dernière boîte, le premier 5 est écrit sur un 3.}

\page{71}
\note{Toute la page est barrée de longs traits obliques.}

$\boxed{(5,5,3)}$
\[
\theta_0 = \theta_1 = \frac{2\pi}{5},\qquad \Theta = \frac{2\pi}{3}
\]
\[
\cos\frac{\theta_0 - \theta_1}{2} = \cos(\theta_0 - \theta_1) = 1,\qquad
\cos\Theta = -1
\]
\[
\sin\frac{\theta_0}{2} = \sin\frac{\theta_1}{2} = \sin\frac{\pi}{5} = 
\]
\note{La valeur qui suit est biffée : \struck{\ill{}}.}
\[
\sin^2\frac{\theta_0}{2} = \frac{1}{2}(1 - \cos\theta_0)
= \frac{1}{2}\Bigl(1 - \frac{\alpha}{2}\Bigr) = \frac{1}{4}(2 - \alpha)
\]
\[
\sin^4\frac{\theta_0}{2} = \frac{1}{16}(\underbrace{\alpha^2 + 4 - 4\alpha}_{5 - 5\alpha})
= \frac{5}{16}(1 - \alpha)
\]
\[
\xi^2 - \frac{1}{\sin^2\frac{\pi}{5}}\,\xi
+ \frac{1}{8}\,\frac{1 - (-1)}{\sin^4\frac{\pi}{5}} = 0
\quad\text{i.e.}\quad
\xi^2 - \frac{1}{\frac{1}{4}(2 - \alpha)}\,\xi
+ \frac{1}{4}\,\frac{1}{\frac{5}{16}(1 - \alpha)} = 0
\]
\[
\xi^2 - \frac{4(2 - \alpha')}{(2 - \alpha)(2 - \alpha')}\,\xi
+ \frac{4}{5}\,\frac{(1 - \alpha')}{(1 - \alpha)(1 - \alpha')}
\]
\[
N(2 - \alpha) = 4 - 1 + 2 = 5,\qquad N(1 - \alpha) = 1 - 1 + 1 = 1
\]
\[
-\alpha' = 1 + \alpha,\qquad 2 - \alpha' = 3 + \alpha,\qquad
1 - \alpha' = 2 + \alpha
\]
\note{Le «~$-1$~» de $\cos\Theta$ est écrit tel qu'il est sur la page, le
chiffre devant $\cos$ biffé. À gauche, un pentagone inscrit dans un cercle
avec un triangle marqué $\zeta$, et la formule
$\frac{1}{|1 + \zeta|}\,\mathcal{R}(1 + \zeta) = \sin\frac{\pi}{5}$.}

\[
\zeta = \frac{\alpha}{2} + i\sqrt{1 - \frac{\alpha^2}{4}}
= \frac{\alpha}{2} + \frac{i}{2}\sqrt{\alpha + 3}
\qquad \Bigl(\sqrt{1 - \tfrac{\alpha^2}{4}} = \tfrac{1}{2}\sqrt{4 - \alpha^2}\Bigr)
\]
$\alpha^2 + \alpha - 1 = 0$, \emph{NB} $4 - \alpha^2 = 3 + \alpha$
\[
|\zeta + 1|^2 = \Bigl(1 + \frac{\alpha}{2}\Bigr)^2
+ \Bigl(\frac{\sqrt{\alpha + 3}}{2}\Bigr)^2
= 1 + \alpha + \frac{\alpha^2}{4} + \frac{\alpha + 3}{4}
= \frac{3}{2} + \alpha
\]
\[
\sin\frac{\pi}{5} = \frac{1 + \alpha/2}{\sqrt{3/2 + \alpha}}
= \frac{1}{2}\,\frac{(\alpha + 2)}{3 + 2\alpha}\sqrt{3/2 + \alpha}
= \alpha\sqrt{3/2 + \alpha}
\]
\note{Suivi d'un terme biffé : \struck{\ill{}}.}
\[
\sin^2\Bigl(\frac{\pi}{5}\Bigr) = \alpha^2(3/2 + \alpha)
= (1 - \alpha)(3/2 + \alpha)
= 3/2 - \tfrac{1}{2}\alpha - \underbrace{\alpha^2}_{\alpha - 1}
= \frac{1}{2}(\alpha + 1)
\]
\[
N(3 + 2\alpha) = 9 - 4 - 6 = -1,\qquad (3 + 2\alpha)(3 + 2\alpha')
\]
\[
\frac{\alpha + 2}{3 + 2\alpha} = -(\alpha + 2)(3 + 2\alpha')
= -(\alpha + 2)(1 - 2\alpha)
= -\bigl[\underbrace{-2\alpha^2}_{2\alpha - 2} - 3\alpha + 2\bigr] = \alpha
\]
\[
\xi^2 - \Bigl[\frac{4}{5}(3 + \alpha)\Bigr]\xi + \Bigl(\frac{4}{5}(2 + \alpha)\Bigr) = 0
\]
\[
\xi = \frac{2}{5}(3 + \alpha) \pm \frac{2}{5}
\sqrt{(3 + \alpha)^2 - 5(2 + \alpha)}
= \frac{2}{5}(3 + \alpha) \pm \sqrt{-5\alpha}
\]
\[
9 + 6\alpha + \alpha^2 - 10 - 5\alpha = \ldots
\qquad (\alpha^2 = 1 - \alpha)
\]
\note{Le résultat, \struck{$5\alpha$}, est biffé.}
\note{Les coefficients sous le radical sont écrits sur d'autres, biffés ;
le dernier résultat est barré. Le calcul est laissé tel quel : avec
$\alpha^2 + \alpha - 1 = 0$, la quantité sous l'accolade est nulle.}

\page{72}
$\boxed{(5,3,5)}$ et $\boxed{(3,5,5)}$
\note{La première boîte est précédée d'un signe de racine.}
\[
\theta_0 = \frac{2\pi}{5},\quad \theta_1 = \frac{2\pi}{3},\quad
\theta_1 - \theta_0 = \frac{4\pi}{15},\quad
\frac{\theta_1 - \theta_0}{2} = \frac{2\pi}{15},\quad
\cos\frac{\theta_1 - \theta_0}{2} = \cos\Theta.
\]
\[
\underbrace{\cos\frac{2\pi}{15}}_{\frac{\alpha}{4} + \frac{1}{4}\sqrt{9 + 3\alpha}}
+ \underbrace{\cos\frac{\pi}{5}}_{\frac{1}{2}(1 + \alpha)}
\overset{?}{\geq}
2\underbrace{\sin\frac{\pi}{5}\sin\frac{\pi}{3}}
_{\frac{\sqrt{3}}{2}\sqrt{3 + \alpha} = \frac{1}{2}\sqrt{9 + 3\alpha}}
\]
\[
\text{i.e.}\quad \frac{1}{2} + \frac{3\alpha}{4} + \frac{1}{4}\sqrt{9 + 3\alpha}
\ \geq\ \ldots
\]
\note{Sous les accolades, des états biffés ($\sqrt{2 + \alpha}$, $1 + 2\alpha$,
$\sqrt{3}$).}

\[
2\cos^2\frac{\pi}{5} = 1 + \cos\frac{2\pi}{5} = 1 + \alpha/2
= \frac{1}{2}(2 + \alpha),\quad
4\cos^2\frac{\pi}{5} = 2 + \alpha = (1 + \alpha)^2,\quad
\cos\frac{\pi}{5} = \frac{1}{2}(1 + \alpha)
\]
\[
\cos\frac{\pi}{5} = \frac{1}{2}\alpha,\qquad
\sin^2\frac{\pi}{5} = 1 - \frac{1}{4}\alpha^2 = \frac{1}{4}(4 - \alpha^2)
= \frac{1}{4}(3 + \alpha)
\]
\struck{$\sin\frac{\pi}{5} = \sqrt{1 - \frac{\alpha^2}{4}}$}
\note{Les deux valeurs de $\cos\frac{\pi}{5}$, $\frac{1}{2}(1 + \alpha)$ et
$\frac{1}{2}\alpha$, sont transcrites telles qu'écrites, l'une sous l'autre.}

$\alpha^2 + \alpha - 1 = 0$
\[
\alpha = \frac{1}{2}(\sqrt{5} - 1) \leq \frac{2}{3}
\]
\struck{car $\sqrt{5} \leq 3$} car $\sqrt{5} \leq \frac{4}{3} + 1 = \frac{7}{3}$
i.e.\ $5 \leq \frac{49}{9}$ i.e.\ $45 \leq 49$~!
\[
\frac{\theta_1 - \theta_0}{2} = \frac{\pi}{5} - \frac{\pi}{3}
\]
\[
\cos(\ ) = \underbrace{\cos\frac{\pi}{5}}_{\alpha/2}\underbrace{\cos\frac{\pi}{3}}_{1/2}
+ \underbrace{\sin\frac{\pi}{5}}_{\frac{1}{2}\sqrt{3 + \alpha}}
\underbrace{\sin\frac{\pi}{3}}_{\sqrt{3}/2}
= \frac{\alpha}{4} + \frac{1}{4}\sqrt{9 + 3\alpha}
\]

\[
\frac{1}{4}\sqrt{9 + 3\alpha} \overset{?}{\leq} \frac{1}{2} + \frac{3\alpha}{4}
\]
\note{Suivi d'un terme biffé : \struck{\ill{}}.}
\[
\sqrt{9 + 3\alpha} \leq 3\alpha + 2
\]
\[
9 + 3\alpha \leq 9\alpha^2 + 4(\ldots) = \ldots
\]
\[
\ldots - \alpha + 8 + 4\alpha = 13 - 3\alpha
\]
\[
6\alpha \leq 4,\qquad \alpha \leq \frac{2}{3}\quad \text{OK}
\]
\note{Termes biffés dans ces lignes : \struck{\ill{}} après $3\alpha + 2$ ;
\struck{\ill{}} dans la parenthèse après $4$ ; \struck{$9 + 8\alpha \leq$}
en tête de la troisième ligne, et \struck{$+ 4\sqrt{2\alpha + \alpha^2}$} ;
\struck{OK} devant $6\alpha \leq 4$.}
\note{Cette colonne de droite est surchargée et en partie hachurée ; seuls
les termes conservés sont donnés. Le résultat final se lit sans doute.}

\struck{Reste $(3,5,5')$}\quad $\boxed{(5,3,5')}$ et $\boxed{(3,5,5')}$

même calcul, sauf qu'ici
$\cos\frac{\Theta}{2} = \cos\frac{\pi}{5} = \frac{1}{2}($\struck{$1 + \alpha$}$)
= -\frac{1}{2}\alpha'$ \ill{} dans le \ill{}, dont on remplace \ill{}
$\cos\frac{2\pi}{5} = -\frac{1}{2}$ \struck{\ill{}}, qui est plus petit, il
faut vérifier
\[
\sqrt{9 + 3\alpha} \leq \frac{\alpha}{\cdot} - \frac{\alpha}{4}\quad \text{pas vrai~!}
\]
\note{Chiffres et termes biffés : un \struck{4} devant la racine, un
terme \struck{\ill{}} après $\leq$, un \struck{4} au dénominateur (noté
$\cdot$).}
\[
9 + 3\alpha \leq 4(1 - \alpha) + \alpha^2 + 4\sqrt{\alpha^2(1 - \alpha)}
\]
\marginal{à étudier (\uncertain{Condition} \ill{} \ill{}~?)}
\note{Note marginale écrite en travers, à gauche de la boîte. Le calcul de
droite est hachuré en plusieurs endroits.}

\page{73}
Cas $\boxed{(5,5,5)}$
\[
1 + \underbrace{\cos\frac{\pi}{5}}_{\frac{1}{2}(1 + \alpha)}
\geq 2\underbrace{\sin^2\frac{\pi}{5}}_{\frac{1}{2}(3 + \alpha)}
\]
\struck{i.e.\ $3 + \alpha$} \ill{} \ill{} OK.
\note{Les chiffres de la boîte sont des 5 écrits sur d'autres.}

Seuls cas indéterminés \ill{} \uncertain{restent} ceux de
\emph{$(5,3,5')$} et \emph{$(3,5,5')$} --- vérifier s'il existe bien ---
sinon, il faut trouver comment \ill{} les \ill{} éliminables\ldots
\note{Le reste de la page est blanc.}


\section{2-Polyèdres réguliers finis en car.\ 0 (classification)}
\note{Titre de sa main sur la couverture (page 74), avec au crayon
«~(6 p)~» : les pages 75 à 80.}

\page{75}
Nous nous intéressons au cas $n = 2$, et au cas où $G_0$ est fini et
\add{a sa représentation dans $\mathrm{Aut}(X)$} \struck{\ill{} une
représentation} fidèle \struck{\ill{} $\mathrm{Aut}(X)$} (\uncertain{quant
à} dire par \ill{} $\mathrm{Aut}(C)$). --- il faut y \ill{} \ill{} \ill{}
\uncertain{trivial} pour $G_0 \to \mathrm{Aut}(C)$]
\marginal{\uncertain{vérifier} \ill{} \ill{}}

\[
\left\lbrace
\begin{array}{ll}
\text{Invariants} & \beta_0 = \alpha_0 + 2,\quad \beta_1 = \alpha_1 + 2 \\
 & \rho = 8 - 2(\beta_0 + \beta_1) = -2(\alpha_0 + \alpha_1),\quad
 \tilde{\rho} = -(\alpha_0 + \alpha_1) \\
 & \sigma = 4 - \beta_1 = 2 - \alpha_1,\quad \sigma' = 4 - \beta_0 = 2 - \alpha_0
\end{array}\right.
\]

$Q_C = Y$ de dim.\ relative 1, lisse ssi $\beta_0\beta_1\tilde{\rho}$
inv., i.e.\ ssi en tout point on a
$\alpha_0 \neq -2$, $\alpha_1 \neq -2$, \struck{\ill{}}
$\alpha_0 + \alpha_1 \neq 0$

$Q_X$ de dim relative 2, lisse ssi $\beta_0\beta_1\sigma$ inv, i.e.\ ssi en
tout point on a $\alpha_0 \neq -2$, $\alpha_1 \neq \pm 2$,

$\rho_f = \tau_0\tau_1$ son ordre dépend de la façon connue de
$\beta_0 = \alpha_0 + 2$

$\rho_s = \tau_1\tau_2$ \ldots\ $\beta_1 = \alpha_1 + 2$,

\struck{\ill{}} On sait que
\[
\mathrm{Aut}(C, Q_C) \xrightarrow{\ \sim\ } \mathrm{Aut}(Q_C) = \mathrm{Aut}(Y)
\]
\add{\uncertain{sauf} si $G_0$}
Donc la cat.\ des 2-polyèdres réguliers \add{\uncertain{pseudo-réflexifs}}
\struck{\ill{}} épinglés sur $S$ \ill{} équivalente à celle des triples
$(Y, \varphi^{Y}_{G_0})$ d'une conique $Y$ sur $S$ et d'un hom
\[
\varphi^{Y}_{G_0} : G_0 \longrightarrow \mathrm{Aut}(Y),
\]
satisfaisant aux conditions
\begin{itemize}
\item[a)] \struck{\ill{}} $\rho_s^2 \neq 1$, $\rho_f^2 \neq 1$ en toute
fibre ($\Rightarrow \tau_0, \tau_1, \tau_2 \neq \mathrm{id}$ sur les fibres,
donc $\tau_0, \tau_1, \tau_2$ des \emph{réflexions})
\item[b)] Les sections $h_0, h_1, h_2$ de $C$ définies par
$\tau_0, \tau_1, \tau_2$) sont non alignées en tout point.
\struck{\ill{}}
\end{itemize}
Si $\beta_0$ inv et $(\beta_1, \cdot)\mathcal{O}_S = \mathcal{O}_S$\note{le second terme,
noté $\cdot$, est surchargé et illisible (\ill{})}
[i.e.\ \struck{\ill{}} \add{en aucun pt on n'a} \struck{\ill{}}
$\alpha_1 = -2$, $\alpha_0 = 2$], alors b) équivaut à
\begin{itemize}
\item[b')] En tout pt, la repr.\ $\varphi^{C}_{G_0}$ \struck{\ill{}} n'a pas
de droite fixe.
\end{itemize}
Si $Y$ lisse, i.e.\ \struck{$\beta_0, \beta_1, \tilde{\rho}$ \ill{}}
\add{en aucun} pt on n'a \struck{\ill{}} \add{\ill{}}
$\alpha_0 = -2$, $\alpha_1 = -2$, $\alpha_0 + \alpha_1 = 0$, alors la
condition b') équivaut à~:
\begin{itemize}
\item[b'')] $Y$ \add{sur aucune fibre} n'admet pas de 2-diviseur stable par
$\varphi^{Y}_{G_0}$.
\end{itemize}
\marginal{\ill{} \ill{} \ill{} \ill{} \ill{} $(\beta_0, 2)\mathcal{O}_S =
\mathcal{O}_S$, \uncertain{i.e.} $\beta_0 \neq 0$ \ill{} car.\ 2 \ldots\
\ill{} a) (\uncertain{avec} b) \ill{} \ill{} \ill{} cas où $\beta_1$,
$\beta_0 = \beta_1$ \ill{} \ill{} \ill{} réflexions \ill{}}
\note{Longue note marginale écrite en biais dans la marge gauche, avec un
grand «~2~» et une flèche vers a) ; une partie en est biffée.}

\page{76}
I) Condition de finitude, sur corps de base $k$.

\emph{Prop.} Si $G_0$ est fini, alors \struck{les} $\beta_0, \beta_1$ sont
algébriques sur le corps premier. L'inverse est vrai si
\struck{$p =$} car.\ $k > 0$.
\marginal{montrer ce \ill{} \ill{} \uncertain{quelconques}}

\struck{\ill{}} Introduisons \struck{les} \add{des solutions} \ill{}
(dans une clôture alg.\ $\bar{k}$ de $k$) $\zeta_i$ ($i = 0, 1$) de
l'équation aux val.\ propres
\[
\zeta_i^2 - \alpha_i\zeta_i + 1 = 0 \qquad \text{i.e.}\qquad
\zeta_i + \zeta_i^{-1} = \alpha_i
\]
Donc pour $i = 0, 1$ fixé, $\zeta_i$ est déterminé mod passage à
$\zeta_i^{-1}$.

Soit $\nu_i$ = ordre de $\zeta_i$, $i = 0, 1$.

Donc $\nu_i$ est premier à $p$, sauf si $\nu_i = 1$ i.e.\ $\zeta_i = 1$,
i.e.\ $\alpha_i = 2$.

\emph{Prop} a) Si $\nu_i$ \struck{$\neq$} \add{$\neq 1, 2$}
\struck{\ill{}} [i.e.\ $\alpha_i \neq \pm 2$ i.e.\
$\beta_i(4 - \beta_i) \neq 0$], alors $\nu_i$ = ordre de $\rho_i$

\emph{NB} ($\rho_0 = \rho_f$, $\rho_1 = \rho_s$) i.e.\ $\alpha_i = -2$,
i.e.\ $\beta_i = 0$ ($\nu_i = 2$, ou $\nu_i = 1$ en car 2)
\note{La parenthèse qui suit «~NB~» est rattachée par un trait souligné
au cas b).}

b) Si \struck{\ill{}} \ill{}, alors $\rho_i$ d'ordre $2p$

c) \struck{Si $\alpha_i = 2 \neq -2$, \ill{}}
Si $\alpha_i = 2 \neq -2$ (i.e.\ $\nu_i = 1$ en car.\ $\neq 2$) alors
$\nu_i = p$
\note{Au bout de la ligne biffée de c), un bloc hachuré où se lisent
«~$\nu_i = p$~» et «~$p \neq 2$~».}

[\emph{NB} Dans b) et c), on convient que $p = \infty$ dans le cas de car.\
nulle]

\emph{Théorème} Supposons $k$ de car.\ 0. Pour que $G_0$ soit fini, il faut
et il suffit que l'on soit dans l'un des \struck{quatre} \add{six} cas
suivants

\struck{1) $\nu_0 = 3$, $\nu_1 = 3$ (tétraèdre) 2) $\nu_0 = 3$, $\nu_1 = 4$
(octaèdre) 2') $\nu_0 = 4$, $\nu_1 = 3$ (cube) 3) $\nu_0 = 3$, $\nu_1 = 5$
(icosaèdre) 3') $\nu_0 = 5$, $\nu_1 = 3$ (dodécaèdre)}
\note{Première liste, barrée de traits obliques.}

\begin{itemize}
\item[I)] $\nu_0 = 3$, $\nu_1 = 3$ (tétraèdre) i.e.\
$\alpha_0 = \alpha_1 = -1$, i.e.\ $\beta_0 = \beta_1 = +1$
\item[II)] $\nu_0 = 3$, $\nu_1 = 4$ (octaèdre) i.e.\ $\alpha_0 = -1$,
$\alpha_1 = 0$ i.e.\ $\beta_0 = 1$, $\beta_1 = 2$
\item[III)] $\nu_0 = 3$, $\nu_1 = 5$ (icosaèdre) i.e.\ $\alpha_0 = -1$,
$\alpha_1$ solution de l'équation
$\boxed{\alpha_1^2 + \alpha_1 - 1 = 0}$
\item[II')] $\nu_0 = 4$, $\nu_1 = 3$ (cube) i.e.\ $\alpha_0 = 0$,
$\alpha_1 = -1$ i.e.\ $\beta_0 = 2$, $\beta_1 = 1$
\item[III')] $\nu_0 = 5$, $\nu_1 = 3$ (dodécaèdre) i.e.\ $\alpha_0$
solution de $\boxed{\alpha_0^2 + \alpha_0 - 1 = 0}$, $\alpha_1 = -1$ i.e.\
$\beta_1 = 1$
\item[IV)] $\nu_0 = \nu_1 = 5$, \emph{et} $\alpha_0 \neq \alpha_1$
(pentagonaux) (i.e.\ $\alpha_0, \alpha_1$ solutions \emph{distinctes} de
$\boxed{\alpha^2 + \alpha - 1 = 0}$
\end{itemize}
\marginal{on \ill{} \ill{} \ill{} \ill{} pythagoriciens \ill{} \ill{}}
\note{Note marginale en travers, en regard de l'accolade qui réunit I à
III et II' ; une autre accolade réunit II' et III'.}

\page{77}
\emph{Remarques} Les cas I, II \add{II'} (correspondant \ill{} à une
structure, définie à isom.\ unique près), \add{qui existe pour tout $k$ de
car.\ 0 (déf.\ sur $\mathbb{Q}$)}, les autres trois correspondant chacun à
deux structures, définies à isom unique près, \struck{celle} qui existent
ssi $k$ contient une solution de $\alpha^2 + \alpha - 1 = 0$, i.e.\ ssi $k$
contient le corps \uncertain{des racines} cycl.\ \ill{} d'ordre 5. Il y a
donc en tout $3 + 3 \times 2 = 9$ types d'isomorphie de «~polyèdres
réguliers finis~» sur corps $k$ de car.\ 0, contenant $\alpha$.
\note{Le 9 est entouré et souligné deux fois.}

\emph{Indications sur la dém.\ du th.} \uncertain{Lisons} OPS \add{$k$ alg.\
clos, et $=$} $k = \mathbb{C}$.

I) Notons d'abord que d'après la prop.\ si ordres $\nu_0, \nu_1$ sont
finis on a \struck{$\beta_0\beta_1$ \ill{}}
$\alpha_0 \neq \pm 2$, $\alpha_1 \neq \pm 2$
\struck{\ill{} $\sigma, \sigma' \neq 0$, donc $Q_X$ lisse}
\struck{(Pour \ill{} avoir $\rho = 0$ i.e.\ $\alpha_0 + \alpha_1 = 0$ i.e.\
\struck{$\alpha_1 = -2\alpha_0$} $\alpha_0 = \alpha$, $\alpha_1 = -\alpha$,
où $\alpha$ est une solution d'une équation \ill{} cyclotomique (définie
par une racine $\nu = \nu_0 \geq 3/2$) \struck{\ill{}
$\sigma = 2 - \alpha_1 = 2 + \alpha \neq 0$} \ill{} serait dans un cas où
\struck{$Q_E = Q_X \otimes Q_C = Q_X Y$ est} \ill{} plus \ill{})}
\note{Bloc encadré et barré de traits obliques.}
les $\alpha_i$ \add{sont} solution d'équations ½-cyclotomiques d'indice
$\nu_0, \nu_1 \geq 3$. Ces solutions sont \emph{réelles}, donc la situation
provient d'une situation sur $k_0 = \mathbb{R}$. Comme \struck{\ill{}}
pour $G_0$ \ill{} fini, $\exists$ forme quadratique invariante définie
positive sur $V$, \struck{\ill{}} $\boxed{\beta\rho \neq 0}$ i.e.\
\emph{\struck{$X$} lisse}.
\note{Devant «~lisse~», une lettre biffée ($X$ ou $Y$).}

$G_0 \subset \mathrm{GP}(1, \mathbb{C})$ groupe fini d'automorphismes de
$Y \simeq \mathbf{P}^1_{\mathbb{C}}$, \struck{\ill{}} soit
$Z = Y/G_0$ \add{$\simeq \mathbf{P}^1_{\mathbb{C}}$}. Utilisons Hurwitz en
tenant que $Y$ est ramifié sur $Z$ en \add{exactement} trois points. Dans le
cas \uncertain{de} 2 pts, on aurait au dessus des \ill{} \ill{} un diviseur
de degré 2 stable sous $G_0$, absurde. Donc il y a trois points de
ramification, et la formule de Hurwitz nous dit que les \struck{\ill{}}
indices de ramification sont $(2, p, q)$ où $p, q$ \add{\uncertain{sont}}
des invariants plus haut, ($(3,3)\,(3,4)\,(3,5)$ \ill{} \ill{} finis), plus
$(2, 2, p)$ (cas diédral) --- mais dans ce dernier cas, il y a au dessus du
\marginal{$\underline{G}_0 = \mathrm{Im}\, G_0$, $G_0 \to \underline{G}_0$
de noyau d'ordre 1 ou \uncertain{2}}
\note{Note marginale encadrée d'une courbe, à gauche.}

\page{78}
pt de ramification d'indice $p$ un diviseur de degré 2 stable par $G_0$, ce
qui est absurde. Donc on est dans l'un des trois cas
$(2,3,3)$, $(2,3,4)$, $(2,3,5)$.

II) \struck{Conditions \ill{}} \add{\ill{} \ill{}} \ill{} montrer que ces
trois cas correspondent \add{au ss-groupes finis de
$\mathrm{GP}(1, \mathbb{C})$} à des classes de conjugaison bien déterminées
de ss-groupes de $\mathrm{GP}(1, k)$~: groupes tétraédral, octaédral,
icosaédral. Pour ceci, considérons les couples $(\nu_0, \nu_1)$
correspondants (cas I, II, III) --- distinguons dans le cas III
$\alpha = \zeta + \zeta^{-1}$, où $\zeta = \exp\frac{2i\pi}{5}$, donc
$\alpha = 2\cos\frac{2\pi}{5}$. \struck{Les \ill{}} Les polyèdres réguliers
correspondants étant \ill{} tétraèdre, octaèdre, icosaèdre --- on constate
aisément, sur les formules explicites, qu'ils forment des réalisations
géométriques \add{strictes} fidèles de ceux-ci. Leurs groupes $G_0$ sont
respectivement \add{\uncertain{ceux}}
\uncertain{$\mathfrak{A}_4$}, $\mathfrak{S}_4 \times \mathbb{Z}/2\mathbb{Z}$,
$\mathfrak{A}_5 \times \mathbb{Z}/2\mathbb{Z}$, et l'image de ceux-ci dans
$\mathrm{GP}(1, k)$ est respectivement $\simeq \mathfrak{A}_4,
\mathfrak{S}_4, \mathfrak{A}_5$. Appelons ceux-ci les groupes pythagoriciens
standard \add{\uncertain{en position}} dans $\mathrm{GP}(1, k)$ (définis
mod conjugaison). Ceci établit l'existence des cas I, II, III, et par
dualité II') III') \struck{l'unicité}. Il reste à prouver que les seuls cas
de p.r.\ réguliers finis qui ne sont pas couverts par les précédents est le
cas IV, et prouver \add{pour ceux-ci} l'existence de IV. Or \add{pour
celle-ci} la construction des pentagônes associés à un icosaèdre est bien
claire. --- cf plus bas la forme que prend cette construction dans notre
contexte actuel.
\note{Les deux premières lettres de groupe, pour le tétraèdre et
l'octaèdre, ont presque la même forme dans sa main ; la première est lue
$\mathfrak{A}_4$ avec doute.}

\struck{III) Revenons au cas de $G_0 \subset \mathrm{GP}$ \ill{} pour
$\tau_0, \tau_1, \tau_2$ fini.}

\page{79}
\note{Le haut de la page (jusqu'à la ligne biffée) est barré de longs
traits obliques et d'un grand arc de cercle, et coché dans la marge.}

III) Notons maintenant que dans les trois cas I II III envisagés, faisant
la construction sur $\mathbb{R}$ et $\mathbb{C}$, on trouve des cartes
cellulaires régulières sur la sphère qui sont \emph{simplement connexes}.
Car on déduit \add{\ill{} de la théorie des cartes} que le groupe
$G_0^{+}$ \add{$(\simeq \mathfrak{A}_4, \mathfrak{S}_4, \mathfrak{A}_5)$}
est engendré par les générateurs $\rho_s, \rho_f$ avec les seules
relations
\[
\rho_s^{\nu_0} = \rho_f^{\nu_1} = (\rho_s\rho_f)^2 = 1.
\]
La théorie \struck{\ill{} II, III,} des groupes fondamentaux de la sphère
moins trois pts, \struck{\ill{}}

\struck{III} montre alors que dans le cas général, le s-groupe fini $G_0$
\struck{\ill{}} dans \ill{} une \emph{quotient} d'un des groupes
$\mathfrak{A}_4, \mathfrak{S}_4, \mathfrak{A}_5$ --- ceci implique déjà par
les ordres de ses éléments \ill{} $\leq 5$, donc
\[
\boxed{3 \leq \nu_0, \nu_1 \leq 5}
\]
\marginal{on \uncertain{pourrait} \ill{} \ill{} \ill{} \ill{}}
\note{Note marginale au crayon, à gauche, sous un petit cercle.}

D'ailleurs, exprimant (sur $k_0 = \mathbb{R}$) que la forme quadratique
\ill{} \ill{} (\ill{} par \ill{} $a_0, a_1, a_2$)
\[
\begin{pmatrix}
2 & -\beta_0 & 0 \\
-\beta_0 & 2\beta_0 & -\beta_0\beta_1 \\
0 & -\beta_0\beta_1 & 2\beta_0\beta_1
\end{pmatrix}
\]
est définie positive, on trouve les conditions
\[
4 - (\beta_0 + \beta_1) > 0 \qquad
\beta_0(4 - \beta_0) > 0
\]
i.e.
\[
\left\lbrace
\begin{array}{l}
\alpha_0 + \alpha_1 < 0 \\
\ldots\quad -2 < \alpha_0, \alpha_1 < +2
\end{array}\right.
\]
\note{En tête des conditions, \struck{$\beta_0, \beta_1 > 0$} ; devant
la seconde ligne de l'accolade, un mot biffé puis un mot illisible
(\struck{\ill{}} \ill{}).}
Conditions néc.\ et suff.\ pour que la forme quadratique soit définie,
pour \struck{\ill{}} ($\alpha_0, \alpha_1$ réels).
\note{Dans la matrice, les coefficients devant $\beta_0$ et $\beta_0\beta_1$
hors de la diagonale sont écrits sur d'autres chiffres, biffés.}

\page{80}
En fait, les relations $-2 < \alpha_0, \alpha_1 < +2$ sont satisfaites
automatiquement, si $\alpha_0, \alpha_1$ solutions d'équations
½-cyclotomiques, correspondant à $\nu_0, \nu_1 \neq 1, 2$ i.e.\
$\nu_0, \nu_1 \geq 3$. Quitte à passer au système dual, on peut donc
supposer, \ill{} les possibilités, que
\[
3 \leq \nu_0 \leq \nu_1 \leq 5
\]
A priori, il y a \struck{$(3,3)\,(3,5)$} $(3,3)$ $(3,4)$ $(3,5)$ (cas I II
III) $(4,4)$, $(4,5)$ et $(5,5)$. Or $(4,4)$ correspond à
\struck{\ill{}} $\alpha_0 = \alpha_1 = 0$, $\alpha_0 + \alpha_1 = 0$,
impossible, et $(4,5)$ à $\alpha_0 = 0$,
$\alpha_1 = \frac{1}{2}(-1 \pm \sqrt{5})$. Mais \emph{OPS}
$\alpha_1 = \frac{1}{2}(-1 + \sqrt{5})$ qui est $> 0$, \struck{\ill{}}
impossible~! Enfin considérons le cas 5,5 OPS
$\alpha_0 = \frac{1}{2}(-1 + \sqrt{5}) > 0$, donc il faut $\alpha_1 < 0$
donc $\alpha_1 = \frac{1}{2}(-1 - \sqrt{5})$ ($\alpha_0 + \alpha_1 = -1$
OK).
\note{Dans les deux «~OPS~», il écrit $\pm$ avec le signe $+$ souligné
deux fois : c'est la détermination retenue.}

On a gagné~!
\note{Le reste de la page est blanc.}

\end{document}
